\documentclass[10pt,a4paper]{article}
\usepackage{amsmath,amssymb,amsthm,graphicx}
\usepackage{xfrac}
% Minimal editorial layout adapter for the pinned JEP source
% sources/jep-2022-195-contact-systolic/pinned-original-native.tex.
% The publisher document class (cedram) is not part of the retained material, and it
% is the class that used to define a handful of plain text abbreviations used by the
% author's prose. This adapter supplies exactly those abbreviations and nothing else:
% no mathematics, no notation, no theorem environment, no counter.
% HoTT/JEP retention convention: CC BY 4.0, attribution to the original authors.
\providecommand{\ie}{i.e.\ }
\providecommand{\eg}{e.g.\ }
\providecommand{\skpt}{\smallskip}
\usepackage[scr=rsfs,cal=euler]{mathalfa}
\let\mathds\mathbb
\let\epsilon\varepsilon
\newcommand{\Changel}{\mathcode`l="7160}
\newcommand{\Changelback}{\mathcode`l="716C}
\newcommand{\NoChangel}[1]{%
\expandafter\let\csname old\string#1\endcsname=#1
\let#1=\relax
\newcommand{#1}{\mathcode`l="716C\csname old\string#1\endcsname\mathcode`l="7160 }%
}

\NoChangel{\log}
\NoChangel{\ln}
\NoChangel{\lim}
\NoChangel{\limsup}
\NoChangel{\liminf}
\NoChangel{\varlimsup}
\NoChangel{\varliminf}
\NoChangel{\varinjlim}
\NoChangel{\varprojlim}

\newenvironment{enumeratei}
{\bgroup\def\theenumi{\roman{enumi}}\def\theenumii{\arabic{enumii}}\begin{enumerate}}
{\end{enumerate}\egroup}
\newcommand{\Psfrac}[2]{(\sfrac{#1}{#2})}
\newcommand{\psfrac}[2]{\sfrac{(#1)}{#2}}
\newcommand{\spfrac}[2]{\sfrac{#1}{(#2)}}
\newcommand\mto{\mathchoice{\longmapsto}{\mapsto}{\mapsto}{\mapsto}}
\newcommand\hto{\mathrel{\lhook\joinrel\to}}

\usepackage{graphicx}

\newtheorem{thm}{Theorem}[section]
\newtheorem{cor}[thm]{Corollary}
\newtheorem{lem}[thm]{Lemma}
\newtheorem{prop}[thm]{Proposition}
\newtheorem{quest}{Question}[section]
\newtheorem{conj}{Conjecture}[section]
\newtheorem{maintheorem}{Theorem}
\theoremstyle{definition}
\newtheorem{rem}[thm]{Remark}
\newtheorem{exm}[thm]{Example}
\newtheorem{ass}[thm]{Assumption}
\renewcommand*{\themaintheorem}{\Alph{maintheorem}}

\DeclareMathOperator{\area}{area}
\DeclareMathOperator{\lcm}{lcm}
\DeclareMathOperator{\const}{const}
\DeclareMathOperator{\Cal}{Cal}
\DeclareMathOperator{\Flux}{\mathrm{Flux}}
\newcommand{\id}{\mathrm{id}}
\newcommand{\Ham}{\mathrm{Ham}}
\newcommand{\std}{\mathrm{std}}
\newcommand{\vol}{\mathrm{vol}}
\newcommand{\fix}{\mathrm{fix}}
\newcommand{\pr}{\mathrm{pr}}
\newcommand{\Tan}{\mathrm{T}}
\newcommand{\N}{\mathds{N}}
\newcommand{\Z}{\mathds{Z}}
\newcommand{\R}{\mathds{R}}
\newcommand{\CC}{\mathds{C}}\let\C\CC
\newcommand{\UU}{\mathcal{U}}
\newcommand{\VV}{\mathcal{V}}

\let\eul e
\let\diff d
\usepackage{cleveref}
\crefname{originaldefinition}{original definition}{original definitions}
\Crefname{originaldefinition}{Original definition}{Original definitions}
\makeatletter
\newlabel{a:transversality}{{3.4}{source}{}{sourcepointer-1}{}}
\newlabel{a:transversality@cref}{{[originaldefinition][0][]3.4}{[1][0][]source}{}{}{}}
\makeatother
\crefname{originallemma}{original lemma}{original lemmas}
\Crefname{originallemma}{Original lemma}{Original lemmas}
\makeatletter
\newlabel{l:nearly_Besse_normalized}{{3.3}{source}{}{sourcepointer-2}{}}
\newlabel{l:nearly_Besse_normalized@cref}{{[originallemma][0][]3.3}{[1][0][]source}{}{}{}}
\makeatother
\title{A Moser diffeomorphism pulling omega back to omega\_0, fixing the boundary}
\author{The Univalent Foundations Program}
\date{doi:10.5802/jep.195; CC BY-SA 3.0 Unported}
\begin{document}
\maketitle
\setcounter{section}{3}
\setcounter{thm}{6}
\noindent\textit{Editorial note.} The counted statement is the author's labelled lemma `l:pullback\_to\_nu0` (source lines 1501-1504); the proof environment starts after one blank line, so the pairing is by adjacency.\par\medskip
\noindent\textit{Editorial note.} The retained context is the author's displayed collar equation for omega\_0, invoked inside the retained proof; it is an equation, so it consumes no theorem number in the retained sequence.\par\medskip
\noindent\textit{Editorial note.} The retained counter values (section 3, theorem offset 6) make the counted lemma print as Lemma 3.7; this article carries no label-number dump, so the number, and the numbers of the two navigation pointers, were inferred by counting the theorem-class environments of section 3 and are reported with that caveat.\par\medskip
\noindent\textit{Editorial note.} The two external references are navigation pointers to the author's transversality assumption and to the normalisation lemma cited for its first two parts; neither is load-bearing in place of the counted proof, which is retained in full.\par\medskip
\noindent\textit{Editorial note.} The retained preamble is a minimal editorial adapter supplying only the plain text abbreviations that the publisher class used to define, followed by the native author macro block of the pinned source.\par\medskip
\noindent\textit{Editorial scope.} One original result of the pinned publisher article, counted once: the Moser diffeomorphism pulling omega back to omega\_0 and fixing the boundary (author label l:pullback\_to\_nu0).\par\medskip
\noindent\textit{Allowed standard prerequisites (editorial).}\begin{itemize}
\item C2-closeness of contact forms. lambda C2-close to lambda\_0 implies nu and omega close in the corresponding norms.. It is the hypothesis and the first remark of the proof at source line 1507.
\item Moser's trick with boundary conditions. One looks for an isotopy fixing the boundary whose pull-backs equal omega\_0, generated by a vector field solving X\_t interior omega\_t = nu\_0 - nu + d f\_t.. It is the framework of the whole proof, source lines 1509-1518.
\item Collar model on the annulus. On a collar A\_C the forms can be written with the functions h\_1, h\_2, h\_3 and the differential identity partial\_r(r h\_2) - partial\_s h\_1 = r(h\_3 - 1).. They are used to choose f and to estimate R\_t, source lines 1525-1577.
\end{itemize}\medskip
\ifdefined\hypertarget\hypertarget{sourcepointer-1}{}\fi
\noindent\textit{Source reference.} Original definition 3.4: sources/jep-2022-195-contact-systolic/pinned-original-native.tex, source lines 1405--1405. Original result invoked by the retained proof. It is a navigation pointer only: the counted statement's own author proof is retained line for line and nothing is replaced by it. The number of this article's own results was inferred by counting the theorem-class environments of its section of the pinned source, which carries no label-number dump; it is not confirmed by the retained bibliography.\par\medskip
\ifdefined\hypertarget\hypertarget{sourcepointer-2}{}\fi
\noindent\textit{Source reference.} Original lemma 3.3: sources/jep-2022-195-contact-systolic/pinned-original-native.tex, source lines 1289--1289. Original result invoked by the retained proof. It is a navigation pointer only: the counted statement's own author proof is retained line for line and nothing is replaced by it. The number of this article's own results was inferred by counting the theorem-class environments of its section of the pinned source, which carries no label-number dump; it is not confirmed by the retained bibliography.\par\medskip
\begin{equation}
\label{e:omega_0_annulus}
\omega_0|_{A_C}
=
\underbrace{\frac{4\pi}{\alpha_1} \big( q_0+\Psfrac{\alpha_1'}{\alpha_1}p_0\big)}_{(*)} r\,\diff r\wedge\diff s.
\end{equation}

\begin{lem}
\label{l:pullback_to_nu0}
If $\lambda$ is $C^2$-close enough to $\lambda_0$, then there exists a diffeomorphism $\kappa:\Sigma\to\Sigma$ such that $\kappa|_{\partial\Sigma}=\id$ and $\kappa^*\omega=\omega_0$. Moreover, $\kappa$ $C^1$-converges to the identity as $\lambda$ $C^2$-converges to $\lambda_0$.
\end{lem}

\begin{proof}
Note that the smallness of $\|\lambda-\lambda_0\|_{C^2}$ implies the smallness of $\|\nu-\nu_0\|_{C^2}$ and $\|\omega-\omega_0\|_{C^1}$. Assumption~\ref{a:transversality} guarantees that $\omega$ is a symplectic form in the interior of $\Sigma$ inducing the same orientation as $\omega_0$. Therefore, the $2$-forms $\omega_t:=t\omega+(1-t)\omega_0$ are symplectic on $\mathrm{int}(\Sigma)$ for every $t\in[0,1]$. They are actually uniformly $C^1$-close to~$\omega_0$ when $\|\lambda-\lambda_0\|_{C^2}$ is small.

We look for an isotopy $\kappa_t:\Sigma\to\Sigma$ such that $\kappa_t|_{\partial\Sigma}\equiv\mathrm{id}$ and $\kappa_t^*\omega_t=\omega_0$. We build the time-dependent vector field $X_t$ realizing such isotopy, \ie $\tfrac{\diff}{\diff t}\kappa_t=X_t\circ\kappa_t$. By differentiating $\kappa_t^*\omega_t$ with respect to $t$, we obtain
\begin{equation}
\label{e:X_Moser_0}
0=\tfrac{\diff}{\diff t}\kappa_t^*\omega_t
=
\kappa_t^* \big( \diff(X\,\lrcorner\,\omega_t) + \omega-\omega_0 \big)
=
\kappa_t^* \diff\big(X\,\lrcorner\,\omega_t + \nu-\nu_0 \big) .
\end{equation}
We define $X_t$ on $\mathrm{int}(\Sigma)$ by the equation
\begin{equation}
\label{e:X_Moser}
X_t\,\lrcorner\,\omega_t = \nu_0 - \nu + \diff f_t
\end{equation}
for a suitable $C^2$-small smooth function $f:\Sigma\times [0,1] \to\R$ to be determined. A~suitable choice of $f$ will guarantee that $X_t$ has a smooth extension to the whole $\Sigma$ with $X_t|_{\partial\Sigma}\equiv0$.

For every connected component $C$ of $\partial\Sigma$, we fix a collar neighborhood $A_C\subset\Sigma$ so that, with the usual suitable identification $[0,\rho)\times S^1$, the differential forms $\omega_0$ can be written as in~\eqref{e:omega_0_annulus}. Actually, up to rescaling the interval $[0,\rho)$, we can even write $\omega_0|_{A_C}$ as
\[
\omega_0|_{A_C}=-r\,\diff r\wedge\diff s,
\]
where $r\in[0,\rho)$ and $s\in S^1$. Lemma~\ref{l:nearly_Besse_normalized}(i-ii) implies $\nu|_C=\nu_0|_C$ and $\omega_z=(\omega_0)_z$ for all $z\in C$. Therefore we can write
\[
(\nu_0-\nu)|_{A_C} = h_1\,\diff r + r\,h_2\,\diff s,
\qquad
\omega|_{A_C} = -r\, h_3\, \diff r\wedge\diff s,
\]
for some smooth functions $h_i:A_C\to\R$. If $\nu_0$ and $\nu$ are $C^2$-close, the function $h_1$ is $C^2$-small, while $h_2$ is $C^1$-small. Moreover, since $\omega_0$ and $\omega$ are $C^1$-close, the function $1-h_3(0,\cdot)$ is $C^0$-small. In particular, up to choosing the annulus $A_C$ small enough, $h_3$ is strictly positive on the whole $A_C$. Since $\diff(\nu_0-\nu)=\omega_0-\omega$, we have
\begin{equation}
\label{e:d_nu_local_coord}
\partial_r(r\, h_2)- \partial_s h_1 = r (h_3 - 1).
\end{equation}
The two-form $\omega_t|_{A_C}$ is given by
\[
\omega_t|_{A_C}=-r\,(t(h_3-1)+1)\,\diff r\wedge\diff s,
\]
and if we write the vector field $X_t|_{A_C}$ in $(r,s)$ coordinates as $X_t=R_t\,\partial_r + S_t\,\partial_s$, Equation~\eqref{e:X_Moser} becomes
\[
R_t = -\frac {r\, h_2 + \partial_s f}{r\,(t(h_3-1)+1)},\qquad
S_t = \frac { h_1 + \partial_r f}{r\,(t(h_3-1)+1)}.
\]
We now choose $f:\Sigma\to\R$ to be a smooth function such that $f|_{\partial\Sigma}\equiv0$ and, on any collar neighborhood $A_C=[0,\rho)\times S^1$ as above, satisfies
\[
f(r,s)
=
\begin{cases}
\displaystyle -\int_0^r h_1(x,s)\,\diff x, & \mbox{if }r\leq\tfrac13\rho, \vspace{5pt} \\
0, & \mbox{if }r\geq\tfrac23\rho.
\end{cases}
\]
We shall choose such an $f$ so that $\|f\|_{C^2}\leq\const\|h_1\|_{C^2}$, and in particular $f$ is $C^2$\nobreakdash-small since $\nu$ and $\nu_0$ are $C^2$-close. With this choice of $f$, we have $S_t(r,s)=0$ if $r\leq\tfrac13\rho$. As for the function $R_t$, for all $r\leq\tfrac13\rho$ Equation~\eqref{e:d_nu_local_coord} implies
\begin{align*}
R_t(r,s)
&=
-\frac {r\, h_2(r,s) + \partial_s f(r,s)}{r\,(t(h_3-1)+1)}
=
-\frac {r\, h_2(r,s) -\int_0^r \partial_s h_1(x,s)\,\diff x}{r\,(t(h_3-1)+1)}\\
&=
-\frac {r\, h_2(r,s) +\int_0^r \big(x (h_3(x,s) - 1) - \partial_x(x\, h_2(x,s))\big)\,\diff x}{r\,(t(h_3-1)+1)}\\
&=
-\frac {\tfrac1r\int_0^r x (h_3(x,s) - 1) \,\diff x}{t(h_3-1)+1}.
\end{align*}
We already know that the function $t(h_3-1)+1$ appearing in the denominator of the above equation is nowhere vanishing. As for the numerator, we can rewrite the integral as
\[
\int_0^r x (h_3(x,s) - 1) \,\diff x = r^2\, h_4(r,s)
\]
for some $C^0$-small smooth function $h_4:A_C\to\R$ such that $h_4(0,s)=h_3(0,s) - 1$. Therefore
\[
R_t(r,s) = -\frac {r\,h_4(r,s)}{t\,(h_3-1)+1}.
\]
From this expression we readily infer that $R_t$ is $C^1$-small, extends smoothly to the whole $\Sigma$, and $R_t|_{\partial\Sigma}\equiv0$. Summing up, we obtained a $C^1$-small smooth vector field~$X_t$ on $\Sigma$ satisfying~\eqref{e:X_Moser} and $X|_{\partial\Sigma}\equiv0$.
Its flow $\kappa_t$ is $C^1$-small for all $t\in[0,1]$, and satisfies $\kappa_t|_{\partial\Sigma}=\id$ and, by~\eqref{e:X_Moser_0}, $\kappa_t^*\omega_t=\omega_0$.
\end{proof}

\end{document}
